\section{Importance of Prompt Precision and Context Design}

The foundations chapter separated prompt precision from context design as two distinct gaps an instruction can leave open: precision concerns whether an instruction states the actual requirement clearly, while context concerns whether the agent has been given the surrounding information --- conventions, existing patterns, project-specific constraints --- it needs to act on that instruction consistently with the rest of the system. The practical project's two parts show what happens when each gap is left unaddressed, and what changes when it is closed deliberately.

Part 1 demonstrated the precision gap directly: output quality tracked how precisely a prompt specified what was wanted, at the level of both visual detail and missing functional requirements, and the resulting gaps were not failures of capability but direct consequences of what the prompt had left unstated. The practical implication is straightforward but easy to underweight in the moment a prompt is written: an agent does not signal that a requirement was ambiguous by producing a visibly uncertain or inconsistent result. It settles on one interpretation and produces output that looks complete and correct, which means the burden of precision falls entirely on whoever writes the prompt, before any output exists to check against. \textbf{A recommendation that follows directly from this is to state assumptions explicitly rather than leaving them implicit --- particularly for requirements a human would consider too obvious to need spelling out, since those are exactly the ones most likely to be silently reinterpreted.}

Part 1 also illustrates the context gap, though less explicitly than the precision gap, since CLAUDE.md began with only a partial set of conventions and grew reactively rather than supplying project context up front. Part 2's response to this was architectural rather than a matter of writing better individual prompts: rather than repeating relevant context inside each task instruction, it maintained a persistent set of project knowledge files split by kind, kept synchronized with the codebase as part of the standard implementation workflow rather than left to go stale. \Cref{lst:frontend-structure-tree} shows a concrete example of this kind of context, kept ready for the agent to consult rather than needing to be rediscovered or restated with each new task. \textbf{The best practice this points to is institutionalizing context design at the level of the project rather than the level of the individual prompt: maintain a standing, up-to-date source an agent can consult for conventions and structure, rather than relying on each task instruction to restate them correctly from scratch.}

Taken together, the two parts suggest that precision and context should be treated as separable practices rather than as one problem solved by writing more detailed prompts. A precise instruction given without adequate context still leaves an agent to invent conventions it has no way of knowing about; thorough context supplied alongside a vague instruction still leaves the actual requirement unresolved. \textbf{Addressing both --- stating the requirement explicitly, and maintaining a persistent source of project context the agent can draw on --- is, on the evidence of both parts, a precondition for the agent's output being checked against what was actually wanted rather than against a plausible-looking guess.}

